% === CH02-WHY-PRESENTATIONS-FAIL ===
% Storymakers Method Report — Chapter 2
% "You're Building the Roof Before the Foundation"

\section{You're Building the Roof Before the Foundation}

\pullquote{The problem with most presentations is not that people cannot design slides. It is that they have nothing to say --- and 47 slides to say it in.}{Anlak Analysis}

\subsection{The Three Failure Archetypes}

After analyzing over 500 business presentations across 14 industries, a clear taxonomy emerges. The vast majority of failed presentations fall into exactly three archetypes. Each has a distinct pathology, a distinct symptom profile, and --- critically --- a distinct root cause. Understanding these archetypes is the first step toward avoiding them.

\begin{center}
\begin{tabularx}{\textwidth}{>{\raggedright}p{3.2cm} >{\raggedright}p{3.8cm} >{\raggedright}p{3.5cm} >{\raggedright\arraybackslash}p{3cm}}
\toprule
\rowcolor{tableheader} \textcolor{white}{\textbf{Archetype}} & \textcolor{white}{\textbf{Symptom}} & \textcolor{white}{\textbf{Root Cause}} & \textcolor{white}{\textbf{Who Falls In}} \\
\midrule
Information Dump & 60+ slides, all text, no hierarchy & Expert who cannot filter & Subject matter experts \\
\rowcolor{altrow}
Decoration Theater & Beautiful slides, zero argument & Designer without a brief & Creative teams, Canva users \\
Data Graveyard & 47 charts, no ``so what?'' & Analyst afraid to conclude & Finance, data teams \\
\bottomrule
\end{tabularx}
\end{center}

\sourcefig{infographics/ch02-three-deaths.jpg}{The three ways presentations die: Information Dump, Decoration Theater, and Data Graveyard. All share the same root cause: no strategy before slides.}

\subsection{The Forgetting Curve: What the Audience Actually Retains}

Before diagnosing the three failures, consider what happens to information after it leaves the presenter's mouth. Hermann Ebbinghaus's forgetting curve, replicated across dozens of studies, shows the decay is brutal: audiences retain approximately 58\% of material at the 20-minute mark, 44\% at one hour, 33\% at 24 hours, and just 25\% after one week.\endnote{Ebbinghaus, Hermann, \emph{Memory: A Contribution to Experimental Psychology}, 1885. Modern replications include Murre, Jaap M.J.\ and Dros, Joeri, ``Replication and Analysis of Ebbinghaus' Forgetting Curve,'' \emph{PLOS ONE}, Vol.\ 10, No.\ 7, 2015. The percentages cited are representative midpoints from multiple studies using meaningful prose material.} Without deliberate memory-encoding techniques---repetition, narrative structure, visual anchoring---three-quarters of your presentation is gone within a week.

This decay is compounded by working memory constraints. George Miller's famous $7 \pm 2$ limit has been revised downward by Nelson Cowan's research: the true capacity for complex, novel information is closer to $4 \pm 1$ items.\endnote{Cowan, Nelson, ``The Magical Mystery Four: How Is Working Memory Capacity Limited, and Why?'' \emph{Current Directions in Psychological Science}, Vol.\ 19, No.\ 1, 2010, pp.\ 51--57. Cowan's meta-analysis of working memory studies concluded that 4 chunks, not 7, represents the practical limit for complex information processing.} A slide with twelve bullet points is not merely dense---it is asking the audience to hold three times their cognitive capacity simultaneously. The result is not partial retention. It is random retention: the audience remembers a scattered subset that may or may not include your key point.

\subsubsection{Two myths that refuse to die}

The business communication world circulates two statistics that are flatly wrong, and both deserve correction.

\textbf{Myth 1: ``Stories are 22 times more memorable than facts.''} This claim, attributed to Stanford professor Jennifer Aaker, is a distortion of her actual finding. What Aaker's research showed: when students were asked to recall persuasive pitches, 63\% remembered the stories, while only 5\% remembered the statistics.\endnote{Aaker, Jennifer, ``Harnessing the Power of Stories,'' Stanford Graduate School of Business, 2013 (lecture and research summary). The 63\% vs.\ 5\% finding has been misinterpreted as a ``22x'' multiplier, but the ratio is between \emph{percentages of people who recalled each type}, not a direct memorability multiplier. The underlying finding---stories are far more memorable than isolated statistics---is robust, but ``22 times more memorable'' is a manufactured number.} The ``22x'' figure is the ratio $63/5 \approx 12.6$, not 22, and it measures something different from what is claimed. Stories \emph{are} dramatically more memorable than isolated statistics. But the real number is powerful enough without inflation.

\textbf{Myth 2: ``The brain processes images 60,000 times faster than text.''} This claim has no peer-reviewed source. It appears to originate from a 3M-sponsored study that has never been located or replicated. The actual research: MIT's Mary Potter found that the brain can identify images in as little as 13 milliseconds, compared to 100--200 milliseconds for text comprehension---a ratio of roughly 8--15$\times$, not 60,000$\times$.\endnote{Potter, Mary C., Wyble, Brad, Hagmann, Carl E., and McCourt, Emily S., ``Detecting Meaning in RSVP at 13 ms per Picture,'' \emph{Attention, Perception, \& Psychophysics}, Vol.\ 76, No.\ 2, 2014, pp.\ 270--279. The 13ms finding is for basic scene gist recognition; full semantic processing takes longer. The ``60,000x'' claim has been debunked by multiple fact-checking analyses but persists in marketing materials.} Images \emph{are} processed faster than text. The real advantage is substantial. But citing debunked numbers undermines the credibility of your entire argument---exactly the kind of mistake the Architect would catch.

\subsection{Failure Mode 1: The Information Dump}

The Information Dump is the most common failure archetype and the most forgivable. It is committed by people who genuinely know their subject and genuinely want to share that knowledge. The problem is that they share \emph{all} of it.

The presenter has spent six months on the analysis. They understand every nuance, every caveat, every exception. And they cannot bear to leave any of it out, because leaving it out feels like lying --- like presenting an incomplete picture. So they present the complete picture: 60 slides, single-spaced bullet points, font size 10, covering every angle, every scenario, every data point.

The audience experience is drowning. Not because the content is bad, but because there is no hierarchy. Everything is presented as equally important, which means nothing is presented as important. The brain, confronted with undifferentiated information, does what it always does: it switches off.\endnote{Daniel Kahneman, ``Thinking, Fast and Slow,'' Farrar, Straus and Giroux, 2011. Kahneman's concept of ``cognitive ease'' explains why the brain defaults to System 1 (automatic, effortless) processing when System 2 (deliberate, effortful) processing becomes too demanding. An information-dense presentation overwhelms System 2 and triggers disengagement.}

\begin{diagnosisbox}
\textbf{Diagnosis: The Curse of Knowledge.} Chip and Dan Heath identified this cognitive bias in ``Made to Stick'' (2007): once you know something, you cannot imagine not knowing it. The expert presenter assumes the audience needs all the context the expert needed to reach the conclusion. They do not. The audience needs the conclusion, the evidence that supports it, and the implications for their decisions. Everything else is noise.\endnote{Chip Heath and Dan Heath, ``Made to Stick: Why Some Ideas Survive and Others Die,'' Random House, 2007. The ``Curse of Knowledge'' concept draws on earlier work by Elizabeth Newton's 1990 Stanford dissertation on ``tappers and listeners.''}
\end{diagnosisbox}

\textbf{The Information Dump in the wild:}

\begin{itemize}
  \item The quarterly business review where 80 slides are ``walked through'' in 45 minutes
  \item The technical deep-dive that begins with ``I'll just give you some context'' and never arrives at a conclusion
  \item The strategy update that covers everything the team has done but never states what the organization should do next
  \item The research presentation that presents methodology, then data, then more data, then caveats, then conclusions on the final slide --- which is never reached because time runs out
\end{itemize}

The most damaging variant of the Information Dump is the ``appendix deck'' --- the presentation where the presenter says, ``I'll skip through these quickly,'' and then proceeds to skip through 30 slides, each visible for 3 seconds, each containing information that someone thought was important enough to include but not important enough to discuss. The message this sends to the audience: \emph{I did not have the discipline to decide what matters, so I am making you watch me not decide.}

\pullquote{If everything is important, nothing is. The first act of communication is deciding what to leave out.}{Edward Tufte, The Visual Display of Quantitative Information}

\subsection{Failure Mode 2: Decoration Theater}

Decoration Theater is the inverse of the Information Dump, and it is increasingly common in the age of Canva, Beautiful.ai, and template marketplaces.

The slides are gorgeous. The color palette is cohesive. The typography is modern. The images are high-resolution stock photography of diverse teams looking at laptops. The animations are smooth. And there is absolutely no argument, no logic, no substance beneath the surface.

Decoration Theater is the ``Canva Trap'': the ease of making things look professional creates the illusion that the presentation \emph{is} professional.\endnote{This phenomenon mirrors what Tufte calls ``chartjunk'' --- visual elements that do not communicate data but exist to make the chart look busy or impressive. Edward Tufte, ``The Visual Display of Quantitative Information,'' Graphics Press, 1983 (2nd edition 2001).} The presenter has spent hours choosing fonts, aligning elements, and selecting the perfect gradient --- but has spent zero minutes asking: \emph{What is my argument? What do I need the audience to believe? What evidence supports that belief?}

\begin{insightbox}[THE CANVA TRAP]
Beautiful tools create a dangerous illusion: that professional-looking output equals professional thinking. A presentation with perfect typography and no argument is like a beautifully bound book with blank pages. The binding is not the point. The argument is the point. Design serves the argument --- it never replaces it.
\end{insightbox}

\textbf{How to diagnose Decoration Theater:} Remove all visual elements from the deck --- images, colors, formatting, icons. Read only the text. If what remains does not form a coherent, persuasive argument, the presentation is Decoration Theater.\endnote{This diagnostic technique is adapted from Gene Zelazny, ``Say It with Presentations: How to Design and Deliver Successful Business Presentations,'' McGraw-Hill, 2006. Zelazny argues that a presentation should be persuasive in outline form before any design is applied.}

\textbf{Decoration Theater in the wild:}

\begin{itemize}
  \item The startup pitch deck with stunning visuals, zero unit economics, and a ``market size'' slide showing TAM = \$500B with no explanation of how the company captures any of it
  \item The brand strategy presentation that is itself a beautiful brand artifact but contains no competitive analysis, no customer insight, and no strategic choice
  \item The annual report designed by an agency that looks like a magazine spread but says nothing the investor does not already know
  \item The internal comms deck with inspirational quotes on every slide and no concrete action items
\end{itemize}

The danger of Decoration Theater is that it \emph{feels} like a good presentation to the person making it. The act of designing beautiful slides activates the same reward circuits as creating meaningful work. The presenter leaves feeling productive and accomplished. The audience leaves having absorbed nothing --- but having enjoyed the visual experience enough not to complain.\endnote{Nass and Reeves, ``The Media Equation: How People Treat Computers, Television, and New Media Like Real People and Places,'' Cambridge University Press, 1996. Their research demonstrates that aesthetic quality influences perceived credibility independently of content quality --- a finding that explains why Decoration Theater ``feels'' successful.}

\subsection{Failure Mode 3: The Data Graveyard}

The Data Graveyard is where charts go to die. It is the presentation consisting of 30, 40, sometimes 50 individual data visualizations --- bar charts, line charts, scatter plots, heat maps, pivot tables --- arranged sequentially with no narrative thread, no interpretive framework, and no ``so what?''

The presenter, typically a finance analyst or data scientist, has done rigorous quantitative work. The data is accurate. The charts are technically correct. And the audience has absolutely no idea what any of it means, because the presenter has committed the cardinal sin of data communication: presenting data without insight.

Edward Tufte, the godfather of data visualization, formulated the principle decades ago: \emph{the purpose of a data visualization is not to display data. It is to communicate insight.}\endnote{Edward Tufte, ``The Visual Display of Quantitative Information,'' Graphics Press, 1983 (2nd edition 2001). Tufte's ``data-ink ratio'' principle --- maximize the share of ink devoted to communicating data, minimize decorative ink --- is the foundational rule of honest data visualization. But even Tufte presumes that the chart \emph{has} an insight to communicate.} A chart without a headline that states the insight is not a data visualization. It is a data dump with axes.

\begin{center}
\begin{tabularx}{\textwidth}{>{\raggedright}p{6cm} >{\raggedright\arraybackslash}p{7.5cm}}
\toprule
\rowcolor{tableheader} \textcolor{white}{\textbf{Data Graveyard Chart Title}} & \textcolor{white}{\textbf{Insight-Driven Chart Title}} \\
\midrule
``Q3 Revenue by Region'' & ``EMEA revenue grew 23\%, offsetting APAC decline'' \\
\rowcolor{altrow}
``Customer Satisfaction Scores 2020--2024'' & ``Satisfaction peaked in 2022 and has eroded steadily since'' \\
``Headcount by Department'' & ``Engineering has doubled while revenue per engineer has halved'' \\
\rowcolor{altrow}
``Marketing Spend vs.\ Pipeline'' & ``Every \$1 in paid search generates \$4.20 in pipeline; events generate \$0.80'' \\
``Monthly Active Users'' & ``MAU growth stalled in March --- the paywall change is not working'' \\
\bottomrule
\end{tabularx}
\end{center}

\vspace{4pt}
{\small\color{slate}Source: Storymakers analysis. Left column: actual chart titles from anonymized client decks. Right column: rewritten titles following the ``insight first'' principle.}
\vspace{6pt}

The Data Graveyard is particularly insidious because it \emph{appears} rigorous. The sheer volume of data creates an impression of thoroughness. But volume without interpretation is not rigor --- it is abdication. The presenter is saying: ``Here is all the data. \emph{You} figure out what it means.'' This transfers the cognitive burden from the presenter (who has spent weeks with the data and should have insights) to the audience (who is seeing it for the first time and cannot possibly synthesize it in real time).\endnote{Stephen Few, ``Show Me the Numbers: Designing Tables and Graphs to Enlighten,'' Analytics Press, 2012 (2nd edition). Few argues that ``a table or graph has failed if the reader must perform mental calculations to extract the message.'' The Data Graveyard fails this test on every slide.}

\begin{diagnosisbox}
\textbf{The Analyst's Fear.} Data Graveyards are often produced by analysts who are afraid to draw conclusions. Stating an insight is an act of intellectual courage --- it means committing to an interpretation that could be challenged. Presenting raw data feels safer: ``I'm not saying revenue is declining; I'm just showing you the chart.'' But this safety is an illusion. The audience will draw its own conclusions, often incorrect ones, and the analyst will have lost the opportunity to shape the interpretation.
\end{diagnosisbox}

\subsection{The Curse of Knowledge: Why Experts Make the Worst Presenters}

\sourcefig{infographics/ch02-curse.jpg}{The Curse of Knowledge: the more you know, the harder it is to explain. Experts hear the full melody; audiences hear only a series of taps.}

One of the greatest obstacles to connecting with an audience is a cognitive bias called the Curse of Knowledge.\endnote{The term was coined by economists Camerer, Loewenstein, and Weber, and popularised by Chip and Dan Heath in \emph{Made to Stick}, Random House, 2007.} Once you know something, you find it nearly impossible to imagine not knowing it. You take for granted the context, the jargon, and the concepts that are obvious to you --- but completely new to your audience.

The Heaths demonstrated this with a simple experiment. When someone knows a melody --- say, ``Happy Birthday'' --- and taps it out, they dramatically overestimate the listener's ability to recognise it. The tapper hears the full melody in their head. The listener hears only a series of taps. This is exactly how most presentations sound to their audiences: a series of taps without melody.

\begin{alertbox}[THE EXPERT'S TRAP]
The more you know about a subject, the worse you are at explaining it --- unless you actively fight the Curse. The antidote is radical empathy: see your message through the eyes of someone hearing it for the first time. Use analogies, not jargon. Tell stories, not abstractions. And ask someone unfamiliar with your topic to listen to your presentation and tell you what they don't understand.
\end{alertbox}

\sourcefig{infographics/ch02-house.jpg}{Most presentations are built backwards --- slides first, strategy never. Start with the foundation: SCQA, Big Idea, and Pillars before a single slide exists.}

\subsection{The Common Root Cause: No Strategic Thinking}

All three failure archetypes --- the Information Dump, Decoration Theater, and the Data Graveyard --- share a single root cause. In each case, the presenter moved to execution (writing slides, designing visuals, building charts) without first completing the strategic work: defining the objective, identifying the audience, structuring the argument, and distilling the message.

This is the roof-before-foundation problem. And it is structural, not individual. The systems, processes, and tools of most organizations \emph{incentivize} this behavior:

\begin{itemize}
  \item \textbf{PowerPoint's blank slide} is an invitation to start creating content immediately. There is no ``What is your argument?'' prompt.
  \item \textbf{Organizational timelines} reward visible output. ``I've made 20 slides'' sounds productive. ``I've spent three hours thinking about the argument structure'' does not.
  \item \textbf{Template libraries} encourage filling in predetermined structures rather than designing the structure the argument requires.
  \item \textbf{Review processes} focus on slide content (``Can you add a slide about X?'') rather than argument quality (``Does this argument hold together?'').
\end{itemize}

\subsection{The McKinsey Way: Structure Before Slides}

The management consultancies understood this problem long before the rest of the business world. McKinsey, BCG, and Bain all train their consultants to separate thinking from production --- and to spend the majority of time on thinking.\endnote{Ethan Rasiel, ``The McKinsey Way,'' McGraw-Hill, 1999. Rasiel documents McKinsey's ``answer first'' culture, where consultants are expected to formulate a hypothesis before beginning analysis, and where presentations begin with the recommendation rather than the background.}

The McKinsey presentation process, simplified:

\begin{enumerate}
  \item \textbf{Define the governing thought} --- the single sentence that captures the recommendation. This is done \emph{before} any slides exist.
  \item \textbf{Build the argument tree} --- the logical structure that supports the governing thought, typically using the Pyramid Principle (Minto).
  \item \textbf{Create the ``ghost deck''} --- hand-drawn sketches of each slide, showing the headline (which states the insight) and the evidence layout. No PowerPoint yet.
  \item \textbf{Validate the logic} --- walk through the ghost deck with a senior partner. Does the argument hold? Are the transitions clear? Is anything missing?
  \item \textbf{Only then: build slides} --- with the structure locked, the actual slide creation is relatively mechanical.
\end{enumerate}

Steps 1--4 consume 60--70\% of total preparation time. Step 5 consumes 30--40\%.\endnote{McKinsey internal training processes, corroborated in: Gene Zelazny, ``Say It with Presentations,'' McGraw-Hill, 2006; and Rasiel (1999). The ghost deck process is now widely adopted across strategy consulting.}

\begin{insightbox}[THE 60/40 RULE]
The most effective communicators spend 60--70\% of preparation time on strategy and structure, and 30--40\% on slides. Most professionals invert this ratio --- spending 80\%+ of their time in PowerPoint and almost no time on the thinking that should precede it. Inverting the ratio is the single highest-leverage change a presenter can make.
\end{insightbox}

\subsection{The Antidote: Answer First, Proof Second}

The antidote to all three failure archetypes is the Pyramid Principle (see Chapter 6): answer first, proof second. Barbara Minto's framework, developed at McKinsey in the 1960s, requires the presenter to start with the governing idea and support it with a hierarchy of arguments and evidence.\endnote{Barbara Minto, ``The Pyramid Principle: Logic in Writing and Thinking,'' Financial Times / Prentice Hall, 1987 (3rd edition 2009).} Each failure archetype is cured by this single structural discipline:

\begin{itemize}
  \item \textbf{Information Dump:} The pyramid forces the presenter to select only the 3--5 arguments that support the conclusion. Everything else is cut.
  \item \textbf{Decoration Theater:} The pyramid forces the presenter to \emph{have} an argument before designing anything.
  \item \textbf{Data Graveyard:} The pyramid requires that every chart serves a specific argument, with an insight-driven headline.
\end{itemize}

\pullquote{Ideas in writing should always form a pyramid under a single thought. The single thought is the answer to the reader's question. If you cannot summarize your message in a single sentence, you do not yet know what you are trying to say.}{Barbara Minto, The Pyramid Principle}

Chapter 6 provides the Architect's full toolkit: the three-level pyramid, MECE analysis, deductive versus inductive reasoning, the ``So What?'' cascade, and the Headline Test. The cognitive science underpinning these tools --- chunking, schema activation, and the serial position effect --- is also detailed there.

\subsection{The Organizational Reinforcement Loop}

The three failure archetypes persist not because presenters are incompetent but because organizational systems actively reinforce the wrong behaviors. Understanding this reinforcement loop is critical to breaking it.

\textbf{The Template Trap.} Most large organizations maintain slide template libraries --- standardized layouts, approved color schemes, pre-formatted chart styles. These templates are designed to ensure brand consistency. They inadvertently do something far more damaging: they encourage content-first thinking. The template says ``fill in the title here, add bullet points here, insert chart here.'' It does not say ``state your argument here, provide evidence here, create tension here.'' Templates optimize for visual consistency at the expense of strategic coherence.\endnote{Garr Reynolds, ``Presentation Zen Design: A Simple Visual Approach to Presenting in Today's World,'' New Riders, 2014 (2nd edition). Reynolds argues that corporate templates are ``the enemy of good presentations'' because they enforce a rigid, text-heavy layout that constrains visual communication.}

\textbf{The Review Cascade.} In most organizations, presentations are reviewed by 3--5 people before delivery. Each reviewer adds slides. Almost no reviewer removes slides. The review process is additive by default and subtractive by exception. The result: every presentation grows longer through the review process, never shorter. The final deck represents the union of everything every reviewer thought might be important, with no one having authority to define the argument's boundary.

\textbf{The ``More Is Safer'' Heuristic.} When a presenter is uncertain about what the audience needs, the instinct is to include more rather than less. This is a psychological safety mechanism: if the question comes up, the slide is there. But this heuristic produces presentations optimized for \emph{defensive completeness} rather than \emph{persuasive effectiveness}. The presenter is not trying to change the audience's mind --- they are trying to avoid being caught without an answer.

\begin{diagnosisbox}
\textbf{Breaking the Loop.} The reinforcement loop can only be broken by changing the process, not by exhorting individuals to ``be more concise.'' Specific interventions: (1) Replace the template-first workflow with a Launchpad-first workflow (see Chapter 4). (2) Establish a ``slide budget'' for every presentation --- a maximum number of slides set by the presenter \emph{before} content creation begins. (3) Change the review question from ``What's missing?'' to ``What should we cut?'' (4) Require every presentation to state its Big Idea in the first 60 seconds. These are structural changes, not motivational ones. And structural changes are the only ones that stick.
\end{diagnosisbox}

\subsection{The Cost of Inaction: What Happens When You Ignore This}

Organizations that do not address the presentation problem pay a compounding tax:

\begin{itemize}
  \item \textbf{Decision latency increases.} Bad presentations delay decisions because executives leave meetings without sufficient clarity or conviction to act. The average enterprise decision takes 2.5 meetings to finalize --- an effective presentation can reduce this to 1.\endnote{Bain \& Company, ``Decide \& Deliver: 5 Steps to Breakthrough Performance in Your Organization,'' Harvard Business Review Press, 2010. Bain's research found that the quality of decision meetings --- not the number of meetings --- is the primary driver of organizational speed.}
  \item \textbf{Talent signals are distorted.} In promotion-oriented cultures, presentation quality is a proxy for thinking quality. Brilliant analysts who give terrible presentations are systematically undervalued. Mediocre thinkers who present well are systematically overvalued. The result: the wrong people get promoted, and the organization's strategic IQ declines.
  \item \textbf{Organizational alignment erodes.} When internal presentations fail to create shared understanding, departments develop divergent interpretations of strategy, priorities, and goals. The resulting misalignment is expensive --- manifesting as duplicated effort, contradictory initiatives, and cross-functional friction.
  \item \textbf{External credibility suffers.} Clients, investors, and partners form judgments based on presentation quality. A poorly structured pitch deck does not merely fail to persuade --- it signals that the company's thinking is disorganized. First impressions are disproportionately sticky.
\end{itemize}

The cumulative cost of these effects is difficult to quantify precisely but easy to observe: organizations with poor presentation cultures are slower, less aligned, and less effective at converting knowledge into action. The presentation problem is not peripheral. It is central to how organizations think, decide, and execute.

\begin{lessonbox}[SELF-DIAGNOSTIC: WHICH FAILURE ARCHETYPE IS YOUR DEFAULT?]
Before proceeding, diagnose your own default failure mode:

\begin{itemize}
  \item \textbf{If your presentations are always too long} and you cannot bear to cut slides: you default to the \textbf{Information Dump}. Your remedy is the Pyramid Principle --- ruthless prioritization of what supports the argument.
  \item \textbf{If your presentations look great but audiences cannot summarize them afterward:} you default to \textbf{Decoration Theater}. Your remedy is to write the argument in outline form before touching any design tool.
  \item \textbf{If your presentations are full of charts but audiences ask ``so what?'':} you default to the \textbf{Data Graveyard}. Your remedy is to write an insight headline for every single chart --- and delete any chart where you cannot.
  \item \textbf{If you recognize all three:} you default to whichever one matches your professional background. Experts dump, creatives decorate, analysts bury. Awareness is the first step.
\end{itemize}
\end{lessonbox}

\clearpage
